\chapter{Bites and Stings from Sea Creatures}

\section*{Definition and Overview}
Bites and stings from sea creatures include stings from jellyfish, bluebottles, stingrays, stonefish, and sea urchins, as well as bites from fish and other aquatic animals. Most are painful but not dangerous; a few, such as those from the box jellyfish and Irukandji syndrome, can be life-threatening.

Australia's coastline is home to several of the world's most venomous marine creatures, so correct first aid and knowing when to call for emergency help are important for anyone swimming, snorkelling, wading, or fishing in coastal waters.

\section*{Epidemiology}
Marine stings are common along tropical and subtropical Australian beaches, especially in summer.

\begin{itemize}
    \item \textbf{Bluebottles:} sting thousands of beachgoers each year and are painful but rarely serious
    \item \textbf{Box jellyfish and Irukandji:} occur mainly in tropical northern waters between November and May, and can be life-threatening
    \item \textbf{Stingrays and stonefish:} less common injuries, usually occurring when a person steps on the creature or handles it while fishing
    \item \textbf{Sea urchins:} frequent cause of puncture wounds in reef and rock-pool areas
    \item \textbf{Most affected groups:} swimmers, snorkellers, and children, particularly during the warmer months
\end{itemize}

\section*{Aetiology and Pathophysiology}
The effects of a marine bite or sting depend on the creature and the toxins involved.

\begin{itemize}
    \item \textbf{Jellyfish tentacles:} discharge microscopic stinging cells into the skin, releasing toxins that cause pain, skin injury, and sometimes systemic effects
    \item \textbf{Bluebottle (Physalia):} venom causes intense local pain and a raised, often linear rash, but rarely affects the heart or breathing
    \item \textbf{Box jellyfish:} potent venom can cause cardiac arrest, breathing failure, and severe skin scarring
    \item \textbf{Irukandji syndrome:} the venom of a tiny jellyfish causes delayed severe pain and a surge in blood pressure, starting hours after what seemed a minor sting
    \item \textbf{Stingrays:} a barbed tail spine causes a puncture wound and can inject venom, causing intense pain and bleeding
    \item \textbf{Stonefish and sea urchins:} spines and venom cause severe pain, swelling, and a risk of retained foreign bodies
    \item \textbf{Anaphylaxis:} some people have severe allergic reactions to marine venoms
\end{itemize}

\section*{Clinical Features}
\begin{itemize}
    \item \textbf{Bluebottle sting:} sharp pain, a raised red welt, often in a line, with itching and swelling at the site
    \item \textbf{Box jellyfish sting:} severe pain, characteristic tentacle marks, and in serious cases collapse, irregular breathing, or cardiac arrest
    \item \textbf{Irukandji syndrome:} initially a minor sting; then severe lower back and abdominal pain, headache, nausea, sweating, and very high blood pressure, typically starting 20 to 40 minutes later
    \item \textbf{Stingray injury:} a deep puncture wound with severe pain and heavy bleeding
    \item \textbf{Stonefish sting:} extreme pain out of proportion to the wound, with swelling and discolouration
    \item \textbf{Sea urchin:} multiple small puncture wounds with embedded spines, pain, and later infection
    \item \textbf{Allergic reaction:} hives, swelling of the lips or throat, wheeze, and dizziness
\end{itemize}

\section*{Differential Diagnosis}
\begin{itemize}
    \item Different jellyfish and marine stings, which require different first aid
    \item An ordinary insect sting or bite
    \item Allergic reaction or anaphylaxis from another cause
    \item Coral or barnacle cuts and scrapes, which frequently become infected
    \item Deep laceration from a fish spine or hook, with or without a retained foreign body
    \item Heat-related illness, drowning, or other beach injuries occurring in the water
\end{itemize}

\section*{When to Seek Medical Care}
See a doctor for any sting that is very painful, covers a large area, or involves a sensitive site such as the eye or face; for puncture wounds from fish spines, stonefish, stingrays, or sea urchins; or for any sting that becomes increasingly red, painful, swollen, or discharging over the following days. Anyone stung by a box jellyfish or a possible Irukandji should be assessed in hospital even if they feel well initially.

\textbf{Call 000 (or local emergency services) for:}
\begin{itemize}
    \item A box jellyfish sting in tropical waters, especially with collapse, breathing difficulty, or confusion
    \item Delayed severe back, chest, or abdominal pain after a minor jellyfish sting, which may indicate Irukandji syndrome
    \item Difficulty breathing, throat swelling, or collapse after any sting
    \item A stingray or stonefish injury with severe bleeding or an embedded barb
\end{itemize}

\section*{Diagnosis and Investigations}
\begin{itemize}
    \item \textbf{History:} description of the creature, timing of the injury, location in the water, and any symptoms
    \item \textbf{Examination:} search for tentacle marks, puncture wounds, and embedded spines, and check blood pressure, pulse, and breathing
    \item \textbf{Blood tests:} including cardiac enzymes and clotting studies in serious envenomations such as the box jellyfish
    \item \textbf{ECG and cardiac monitoring:} in suspected box jellyfish envenomation
    \item \textbf{Imaging:} X-ray or ultrasound to locate embedded fish spines, stingray barbs, or sea urchin fragments
    \item \textbf{Wound swabs:} for infection if the wound becomes infected later
\end{itemize}

\section*{Management}
\begin{itemize}
    \item \textbf{Bluebottle and most other jellyfish stings:} rinse with seawater, remove any visible tentacles with gloved hands, and apply heat or a cold pack for pain
    \item \textbf{Tropical jellyfish (box jellyfish and Irukandji):} pour vinegar over the tentacles, call for help, and arrange hospital treatment; hot water is not used for these stings
    \item \textbf{Stonefish:} soak the limb in hot (not scalding) water for pain relief, then seek hospital care
    \item \textbf{Stingray injury:} control bleeding with direct pressure, keep the wound clean, and seek urgent care, because the barb may need surgical removal
    \item \textbf{Sea urchin spines:} soak in hot water for pain, then seek medical removal of retained spines
    \item \textbf{Antivenom:} available in hospital for box jellyfish and stonefish stings
    \item \textbf{Pain relief:} simple analgesics initially, with stronger measures in hospital for severe pain
    \item \textbf{Wound care and antibiotics:} for puncture wounds at high risk of infection
    \item \textbf{Observation:} in hospital for several hours whenever Irukandji syndrome or a serious envenomation is possible
\end{itemize}

\section*{Complications}
\begin{itemize}
    \item Cardiac arrest and death from severe box jellyfish envenomation
    \item Complications of Irukandji syndrome from very high blood pressure, including stroke and heart damage in rare cases
    \item Scarring and persistent skin discolouration from jellyfish tentacle injuries
    \item Infection of puncture wounds, including severe tissue infection
    \item Retained spines or barbs causing chronic pain or foreign-body infection
    \item Secondary drowning or further injury from collapsing in the water
\end{itemize}

\section*{Prognosis and Outcomes}
Most bluebottle and common jellyfish stings settle within hours to days with simple treatment and leave no lasting harm. With rapid emergency care and antivenom, even box jellyfish and Irukandji victims usually recover fully, although serious cases can be fatal or leave significant scarring.

Puncture wounds from fish spines and stingrays heal well once cleaned and any foreign body is removed, but can become infected if neglected. Full recovery from a severe envenomation may take several weeks.

\section*{Prevention and Self-Care}
\begin{itemize}
    \item Swim at patrolled beaches and follow local advice about jellyfish warnings and stinger nets
    \item In tropical waters, wear a full-body stinger suit during jellyfish season
    \item Shuffle your feet when wading in shallow water to avoid stepping on stingrays
    \item Do not touch or handle marine creatures, including dead jellyfish washed ashore
    \item Wear protective footwear on coral and rocky reefs
    \item Learn the correct first aid for the local area before swimming
    \item Seek medical care promptly for any deep puncture wound from a marine creature
\end{itemize}

\section*{Sources and Further Reading}
The following organisations provide reliable information on marine bites and stings for patients, families, and health professionals.

\begin{itemize}
    \item NHS. \textit{Information on jellyfish stings and their treatment.} https://www.nhs.uk/
    \item Mayo Clinic. \textit{Overview of jellyfish stings and first aid.} https://www.mayoclinic.org/
    \item World Health Organization. \textit{Resources on venomous marine injuries and first aid.} https://www.who.int/
    \item MedlinePlus. \textit{Health information on marine stings and bites.} https://medlineplus.gov/
    \item Centers for Disease Control and Prevention. \textit{Resources on venomous marine animals and injuries.} https://www.cdc.gov/
\end{itemize}
